% BHU PYQ -- Manually transcribed & error-corrected
% Paper: BCF-213 : Corporate Taxation
% Exam: B.Com. (Hons.) FMM Semester III Examination, 2013-14

\documentclass[11pt,a4paper]{article}
\usepackage[a4paper,top=2.5cm,bottom=2.5cm,left=2.8cm,right=2.8cm,headheight=14pt]{geometry}
\usepackage{amsmath,amssymb}
\usepackage[T1]{fontenc}
\usepackage[utf8]{inputenc}
\usepackage{lmodern,microtype,fancyhdr,enumitem,hyperref,lastpage,parskip}

\hypersetup{
    colorlinks=true,
    urlcolor=black,
    linkcolor=black,
    pdftitle={BCF-213: Corporate Taxation -- BHU B.Com. (Hons.) FMM Sem III 2013-14}
}

\pagestyle{fancy}
\fancyhf{}
\renewcommand{\headrulewidth}{0.4pt}
\renewcommand{\footrulewidth}{0.4pt}
\fancyhead[L]{\small   \textit{Banaras Hindu University}}
\fancyhead[R]{\small   \textit{BCF-213: Corporate Taxation}}
\fancyfoot[C]{\small Page  \thepage\ of \pageref{LastPage}}


\newlist{parts}{enumerate}{2}
\setlist[parts,1]{label=(\alph*),leftmargin=*,itemsep=5pt,topsep=3pt}
\setlist[parts,2]{label=(\roman*),leftmargin=*,itemsep=3pt,topsep=2pt}

\newcommand{\pts}[1]{\hfill   \textit{[#1]}}


\begin{document}


\begin{center}
  {\large\bfseries BANARAS HINDU UNIVERSITY}\\[2pt]
  {
\normalsize B.Com. (Hons.) FMM --- Semester III Examination, 2013-14}\\[6pt]
  \rule{0.8\linewidth}{0.5pt}\\[6pt]
  {\large\bfseries Paper No. BCF-213: Corporate Taxation}\\[4pt]
  {
\normalsize Time: 3 Hours\hfill Maximum Marks: 70}
\end{center}

\rule{\linewidth}{0.4pt}

\smallskip



\noindent   \textit{Instructions: Answer all questions. All questions carry equal marks (14 marks each).}

\bigskip




\noindent   \textbf{Question 1.} \\
Write notes on the following:\pts{4.5+4.5+5}

\begin{parts}
  \item Assessee
  \item Assessment Year
  \item Net Taxable Income
\end{parts}

\centerline{   \textbf{OR}}

Explain the characteristics of income tax.\pts{14}

\medskip\rule{\linewidth}{0.2pt}




\noindent   \textbf{Question 2.} \\
State the conditions which an individual must fulfil in order to be called:\pts{5+5+4}

\begin{parts}
  \item Resident
  \item Not-ordinarily resident
  \item Non-resident
\end{parts}

\centerline{   \textbf{OR}}

From the following Profit \& Loss A/c of Mr. Rajkumar for the year ended on 31st March, 2013, compute Income from Business/Profession:\pts{14}

\begin{center}
\renewcommand{\arraystretch}{1.2}

\begin{tabular}{|l|r|l|r|}
\hline
   \textbf{Expenses} &   \textbf{Amount (Rs.)} &   \textbf{Income} &   \textbf{Amount (Rs.)} \\ \hline
General expenses & 5,000 & Gross profit & 1,10,000 \\
Bad debts & 2,500 & Commission & 10,000 \\
Insurance & 750 & Brokerage & 11,000 \\
Provision for bad debts & 1,500 & Sundry receipts & 4,000 \\
Staff salary & 20,000 & Bad debts recovered (earlier allowed) & 2,500 \\
Salary of Rajkumar & 5,000 & Interest on debentures & 5,000 \\
Interest on overdraft & 5,000 & Interest on deposits & 3,500 \\
Interest on loan from Mrs. Rajkumar & 2,500 & & \\
Interest on capital of Mr. Rajkumar & 3,500 & & \\
Depreciation on machinery & 8,000 & & \\
Advertisement expenses & 7,500 & & \\
Contribution to EPF & 5,000 & & \\
Net profit & 79,750 & & \\ \hline
   \textbf{Total} &   \textbf{1,46,000} &   \textbf{Total} &   \textbf{1,46,000} \\ \hline
\end{tabular}
\end{center}

Other informations are as follows:

\begin{enumerate}[label=(\roman*),itemsep=2pt,topsep=2pt]
  \item Depreciation allowed on machinery is Rs. 6,500.
  \item Advertisement expenses include Rs. 4,000 spent on a neon sign board affixed on the office premises.
  \item An income of Rs. 750 accrued during the previous year is not recorded in P \& L A/c.
  \item General expenses include Rs. 1,100 spent on the party of a friend.
\end{enumerate}

\medskip\rule{\linewidth}{0.2pt}




\noindent   \textbf{Question 3.} \\
A, B and C are equal partners in a trading firm. As per partnership deed A and B who are working partners are entitled to a salary of Rs. 40,000 each. Interest payable to all partners @ 12\% per annum amounted to Rs. 4,000, Rs. 5,000 and Rs. 7,000 respectively. C is entitled to receive rent of Rs. 4,000 p.m. of the building owned by him and let out to the firm. In addition to salary B is entitled to a bonus of Rs. 20,000. All the partners are allowed a commission on sale, affected by each of them, which amounted to Rs. 10,000, Rs. 15,000 and Rs. 5,000 respectively. The amount of depreciation debited to P \& L A/c is Rs. 30,000 and the allowable u/s 32 is Rs. 34,000. The amount of entertainment expenditure debited is Rs. 21,000. Net profit as per P \& L A/c after debiting the interest, salary, bonus, commission, etc., is Rs. 50,000.
\smallskip
Compute the book profit and taxable income of the firm for the assessment year 2013-14.\pts{14}

\centerline{   \textbf{OR}}

Explain the salient features of company assessment. Also discuss deductions allowed in computing total income of a company.\pts{14}

\medskip\rule{\linewidth}{0.2pt}




\noindent   \textbf{Question 4.} \\
Sri Ram is getting a salary of Rs. 80,000 per month from a concern. During the financial year, he contributed Rs. 20,000 to Recognised Provident Fund, Rs. 40,000 to Public Provident Fund, Rs. 1,000 per month as Life Insurance Premium and Rs. 15,000 to the National Defence Fund.
\smallskip
Compute the amount of tax to be deducted at source per month during previous year for the assessment year 2013-14.\pts{14}

\centerline{   \textbf{OR}}

When does the liability for advance tax arise? On what date and to what extent it is paid?\pts{14}

\medskip\rule{\linewidth}{0.2pt}




\noindent   \textbf{Question 5.} \\
What do you understand by Best Judgement Assessment? Discuss the provisions related to it under Income-tax Act, 1961.\pts{14}

\centerline{   \textbf{OR}}

Give the names of various Income-tax Authorities. Also discuss the powers of Income-tax Officer.\pts{14}

\vfill
\rule{\linewidth}{0.4pt}

\begin{center}\small
\href{https://bhu-exam.vercel.app/}{Click here to visit our website}\\[4pt]
\href{https://me-aryan.vercel.app/}{Click here to visit our contributor}\\[4pt]
\href{https://quashan-me.vercel.app/}{Click here to visit our contributor}
\end{center}

\end{document}
